% to reference the datset: \cite{zhang2015character}

% Dataset info (where, size, type, classes, etc.)

Dataset for this text classification by topic task is the textual AG News dataset from Hugging Face \cite{zhang2015character}. Each data point has an associated label, the news title and its description. There are 4 classes, each class has 30000 training and 1900 testing samples: world, sports, business, and sci/tech.

% Split details
% split sizes
% split seed

We use the official Hugging Face split (120,000 train and 7,600 test). We split the training set into a new training set (of size 108,000) and a development set (of size 12,000) using stratified sampling with a fixed random seed of 67 to ensure reproducibility. The final split sizes are: train 108,000, dev 12,000, and test 7,600.

\subsection{Preprocessing}
% Preprocessing (Normalization & Tokenization)
% Lowercasing
% Removing whitespace
% Removing punctuation
We combined the title and description of each news article into a single text input, which was then preprocessed. First we lowercased the text, removed extra whitespace, and removed punctuation. We also removed English stopwords using the NLTK library, as they often carry less information and can introduce noise into the model. Finally, we applied lemmatization using NLTK's WordNetLemmatizer to reduce words to their dictionary forms, which helps to reduce the dimensionality of the feature space and improve model performance.

The reason why we remove stop words, is because common words often carry less information and carry noise into our classifier.

% Lemazation (vs. stemming)
The last step in the normalization is lemmatization, the act of reducing words into their dictionary form (e.g. "running" to "run"). Using the NLTK's WordNetLemmatizer, we reduced each word in the text to their dictionary forms to reduce the computational complexity. In this way, different forms of the same word are reduced to a single form.

% Vectorization through IT-IDF

\subsection{Tokenization}

Even though there is no explicit tokenization in our code, texts are vectorized with TF-IDF, which implicitly tokenizes the text into unigrams and bigrams. The TF-IDF vectorizer is configured to consider both unigrams and bigrams (n-gram range of (1, 2)), which allows it to capture more contextual information than unigrams alone while still being less sparse than bigrams alone. We also set a maximum feature limit of 10,000 and a minimum document frequency of 2 to further reduce noise and focus on more informative features.

\subsection{Reproducibility}

All experiments make use of the fixed random seed 67.  We use it for the data split and model initialization where it is applicable.
